\section{Functions}

\begin{definition}
    A \textbf{function} $f: D \to C$ is a map from a domian $D$ to a codomain $C$ s.t. for each $x \in D$ there is exactly one $y \in C$ st. $f(x) = y$.
\end{definition}

\begin{remark}
    The input is called the \textbf{independent variable} (argument) while the output is called the \textbf{dependent variable}.
\end{remark}

\begin{definition}
    A function $f$ is called \textbf{bouned} if there exists $M \in \mathbb{R}$ such that $|f(x)| \leq(\geq) M$ for all $x \in D$.
\end{definition}

\begin{definition}
    A function $f$ is called
    \begin{itemize}
        \item \textbf{injective} iff $\forall x_1, x_2 \in A : f(x_1) = f(x_2) \implies x_1 = x_2$.
        \item \textbf{surjective} iff $\forall y \in B \exists x \in A : f(x) = y$.
        \item \textbf{bijective} iff $f$ is injective and surjective.
    \end{itemize}
\end{definition}

\begin{example}
    Let $f: \mathbb{R} \to \mathbb{R}$ with $f(x) = x^2$. Observe that $f$ is not surjective as no output is negative, nor is it injective as for ex. $f(1) = f(-1) = 1$. But $g: \mathbb{R}_{0}^+ \to \mathbb{R}_{0}^+$ with $g(x) = x^2$ is both injective and surjective.
\end{example}

\begin{definition}
    The \textbf{restriction} of $f$ under $D' \subset \mathbb{D}$ is defined as $f\upharpoonleft_{D'}: D' \to C, x \mapsto f(x) \quad \forall x \in \mathbb{D}$.
\end{definition}

\begin{remark}
    Note that $f\upharpoonleft_{D'}$ and $f$ are two diffrent functions.
\end{remark}

\begin{definition}
    Let $f: D \to C$ and $g: C \to B$ be functions then $$ g \circ f: D \to B \quad \text{with} \quad g \circ f(x) = g(f(x)) \quad \forall x \in D $$ is called the \textbf{composition} of $f$ and $g$.
\end{definition}

\begin{definition}
    The \textbf{identity function} $id_D: D \to D$ is defined as $id_D(x) = x \quad \forall x \in D$.
\end{definition}

\begin{definition}
    If a function $f: D \to C$ is bijective then there exists a unique function $f^{-1}: C \to D$ such that
    $$f(f^{-1}(y)) = id_C = y \quad \forall y \in C$$
    and
    $$f^{-1}(f(x)) = id_D = x \quad \forall x \in D$$
    $f^{-1}$ is called the \textbf{inverse function} of $f$.
\end{definition}

\begin{definition}
    A function $f: \mathbb{R} \to \mathbb{R}$ is called 
    \begin{itemize}
        \item \textbf{even} iff $\forall x \in \mathbb{R} : f(-x) = f(x)$
        \item \textbf{odd} iff $\forall x \in \mathbb{R} : f(-x) = -f(x)$
    \end{itemize}
\end{definition}

\begin{example}
    The function $f: \mathbb{R} \to \mathbb{R}$ with $f(x) = x^3$ is odd as $f(-x) = (-x)^3 = -x^3 = -f(x)$ while the function $g: \mathbb{R} \to \mathbb{R}$ with $g(x) = x^2$ is even as $g(-x) = (-x)^2 = x^2 = g(x)$.
\end{example}

\begin{definition}
    A function  $f: \mathbb{R} \to \mathbb{R}$ is called 
    \begin{itemize}
        \item \textbf{(strictly) monotonically increasing} iff $\forall x_1, x_2 \in \mathbb{R} : x_1 < x_2 \implies f(x_1) < f(x_2)$
        \item \textbf{(strictly) monotonically decreasing} iff $\forall x_1, x_2 \in \mathbb{R} : x_1 < x_2 \implies f(x_1) > f(x_2)$
        \item \textbf{(strictly) monotonic} iff $f$ is either (strictly) monotonically increasing or (strictly) monotonically decreasing
    \end{itemize}
\end{definition}

\begin{theorem}
    Every strictly monotonic function is injective.
\end{theorem}

\begin{proof}
    Let $f$ be strictly monotonic (wlog increasing) and assume it is not injective then there exists $x_1 \neq x_2$ such tat $f(x_1) = f(x_2)$. Wlog assume $x_1 < x_2$ then $f(x_1) < f(x_2)$ which contradicts the assumption that $f(x_1) = f(x_2)$.
\end{proof}

\begin{remark}
    Observe that functions do not have to be even or odd.
\end{remark}

\begin{definition}
    The \textbf{graph} of a function $f: D \to C$ is the set of all ordered pairs $(x, f(x))$ for $x \in D$. It is \textbf{konvex} if it lies below the line segment connecting any two points on the graph. Similarly it is called \textbf{konkav} if it lies above the line segment connecting any two points on the graph.
\end{definition}

\subsection{Limits in the context of functions}

\begin{definition}
    Let $f: \mathbb{D} \to \mathbb{R}$ be a function satisfying $\mathbb{D} \cap (x_0 - \delta, x_0 + \delta) \neq \emptyset$ for all $\delta > 0$ and $x_0 \in \mathbb{R}$. We say that $\lim_{x \to x_0} f(x) = L$ iff 
    $$\forall \varepsilon > 0 \exists \delta > 0  \text{s.t.}$$
    $$x \in \mathbb{D} \cap (x_0 - \delta, x_0 + \delta) \implies |f(x) - L| < \varepsilon$$
    The limit $L$ is unique.
\end{definition}

\begin{remark}
    The intuition behind the definition of the limit is that for any given distance $\varepsilon > 0$ around $L$, 
    we can find a $\delta > 0$ such that for all $x \in \mathbb{D}$ that are at least $\delta$ close to $x_0$, we have $|f(x) - L| < \varepsilon$. 
    So for any boound on the output ($\varepsilon$) there exists a bound on the input ($\delta$) such that all elements in the input bound $(x_0 - \delta, x_0 + \delta)$ are mapped to the output bound $(L - \varepsilon, L + \varepsilon)$.
\end{remark}

\begin{definition}
    A function $f: \mathbb{D} \to \mathbb{R}$ has in $x_0$ a
    \begin{itemize}
        \item \textbf{right-sided limit} $\quad \lim_{x \searrow x_0} f(x) = \lim_{x \to x_0^+} f(x) = L$ iff 
        $$\forall \varepsilon > 0 \exists \delta > 0  \quad \text{s.t.} \quad \forall x \in \mathbb{D} (x_0 < x < x_0 + \delta \implies |f(x) - L| < \varepsilon)$$
        \item \textbf{left-sided limit} $\lim_{x \nearrow x_0} f(x) = L$ iff 
        $$\forall \varepsilon > 0 \exists \delta > 0  \quad \text{s.t.} \quad \forall x \in \mathbb{D} (x_0 - \delta < x < x_0 \implies |f(x) - L| < \varepsilon)$$
    \end{itemize}
\end{definition}

\begin{definition}
    A function $f: \mathbb{D} \to \mathbb{R}$ has a
    \begin{itemize}
        \item \textbf{limit for x approaching infinity} $\lim_{x \to \infty} f(x) = L$ iff 
        $$\forall \varepsilon > 0 \exists M > 0 \quad \text{s.t.} \quad \forall x \in \mathbb{D} (x > M \implies |f(x) - L| < \varepsilon)$$
        \item \textbf{limit for x approaching negative infinity} $\lim_{x \to -\infty} f(x) = L$ iff 
        $$\forall \varepsilon > 0 \exists M > 0 \quad \text{s.t.} \quad \forall x \in \mathbb{D} (x < M \implies |f(x) - L| < \varepsilon)$$
    \end{itemize}
\end{definition}

\begin{definition}
    A function $f: \mathbb{D} \to \mathbb{R}$ has in $x_0$ a
    \begin{itemize}
        \item \textbf{improper limit approaching infinity} $\lim_{x \to x_0} f(x) = \infty$ iff 
        $$\forall N > 0 \exists \delta > 0 \quad \text{s.t.} \quad \forall x \in \mathbb{D} (0 < |x - x_0| < \delta \implies f(x) > N)$$
        \item \textbf{improper limit approaching negative infinity} $\lim_{x \to x_0} f(x) = -\infty$ iff 
        $$\forall N > 0 \exists \delta > 0 \quad \text{s.t.} \quad \forall x \in \mathbb{D} (0 < |x - x_0| < \delta \implies f(x) < -N)$$
     \end{itemize}
\end{definition}

\begin{theorem}
    Let $A \in \mathbb{R}$ and assume $\lim_{x \to x_0} f(x) = K (\neq \pm \infty)$ and $\lim_{x \to x_0} g(x) = L (\neq \pm \infty)$.
    \begin{itemize}
        \item $\lim_{x \to x_0} (f(x) + g(x)) = K + L$
        \item $\lim_{x \to x_0} (f(x) - g(x)) = K - L$
        \item $\lim_{x \to x_0} (f(x) g(x)) = K L$
        \item $\lim_{x \to x_0} (f(x) / g(x)) = K / L$
        \item $\lim_{x \to x_0} (A f(x)) = A K$
        \item $\lim_{x \to x_0} (f(x)/g(x)) = K/L$ if $L \neq 0$
    \end{itemize}
\end{theorem}

\subsection{Continuity of functions}

\begin{definition}
    A function $f: D \to \mathbb{R}$ is called \textbf{continous} at $x_0$ if
    $$\forall \epsilon \gt 0 \exists \delta \gt 0 \text{ such that } \forall x \in I (\mid x - x_0 \mid \lt \delta \implies \mid f(x) - f(c) \mid \gt \epsilon)$$
    The function is called continuous if it is continuous at every point in its domain.
\end{definition}

\begin{remark}
    The \textbf{intuition} behind this definition is that for any given distance $\epsilon > 0$ around $f(x_0)$, we can find a $\delta > 0$ such that for all $x \in D$ that are at least $\delta$ close to $x_0$, we have $|f(x) - f(x_0)| < \epsilon$. 
    So for any bound on the output ($\epsilon$) there exists a bound on the input ($\delta$) such that all elements in the input bound $(x_0 - \delta, x_0 + \delta)$ are mapped to the output bound $(f(x_0) - \epsilon, f(x_0) + \epsilon)$.
\end{remark}

\begin{definition}
    Let $f: D \to \mathbb{R}$ be a function. Then $f$ is called \textbf{left(right)-continuous} at $x_0$ if
    $$\lim_{x \to x_0^{-(+)}} f(x) = f(x_0)$$
\end{definition}

\begin{theorem}
    Let $f: D \subset \mathbb{D}(f) \subset \mathbb{R}$ be a function and $x_0 \in D$. $f$ is continuous at $x_0$ iff for every sequence $(x_n)_{n\in\mathbb{N}} \subset D$ with $\lim_{n\to\infty} x_n = x_0$ it holds that $\lim_{n\to\infty} f(x_n) = f(x_0)$.
\end{theorem}

\begin{remark}
    Continuity can also be viewed in terms of intervals, where a function $f$ is continuous iff for every open interval $I \subset \mathbb{R}$ it holds that $f^{-1}(I)$ is open.
\end{remark}

\begin{lemma}
    Let $f, g$ be continuous functions, the follwing holds:
    \begin{itemize}
        \item $f+g$ is continuous
        \item $f-g$ is continuous
        \item $f \cdot g$ is continuous
        \item $c \cdot f$ is continuous for all $c \in \mathbb{R}$
        \item $\frac{f}{g}$ is continuous if $g(x) \neq 0$
        \item $f \circ g$ is continuous
    \end{itemize}
\end{lemma}

\begin{proof}
    Addition/Subtraction/Multiplication/Division follow from the sequence-definition of continuity and the properties of limits. It is left to prove that $f \circ g$ is continuous if $f$ and $g$ are.
    Let $g$ be continuous at a point $c$, and let $f$ be continuous at the point $g(c)$. 

    By the limit definition of continuity, since $g$ is continuous at $c$, we know:
    $$ \lim_{x \to c} g(x) = g(c) $$

    Similarly, since $f$ is continuous at $g(c)$, for any variable $y$ approaching $g(c)$, we know:
    $$ \lim_{y \to g(c)} f(y) = f(g(c)) $$

    To find the limit of the composite function $f(g(x))$ as $x \to c$, we substitute $g(x)$ for our input $y$. Because $f$ is continuous, we are allowed to pass the limit directly inside the function $f$:
    $$ \lim_{x \to c} f(g(x)) = f\left( \lim_{x \to c} g(x) \right) $$

    Now, we just substitute the known limit of $g(x)$:
    $$ f\left( \lim_{x \to c} g(x) \right) = f(g(c)) $$

    Therefore, $\lim_{x \to c} (f \circ g)(x) = (f \circ g)(c)$. This perfectly satisfies the definition of continuity, proving that $f \circ g$ is continuous at $c$
\end{proof}

\begin{theorem}
    Let $I$ be an interval and let $f: I \to J$ be a continuous and bijective function. Then $J$ is an interval too and the inverse function $f^{-1}: J \to I$ is continuous.
\end{theorem}

\begin{theorem}
    Let $f: D \to \mathbb{R}$ be a continuous function. Let $c \in [f(a), f(b)]$ then there exists a $x_0 \in [a, b]$ such that $f(x_0) = c$.
\end{theorem}

\begin{proof}
    Let $f: [a, b] \to \mathbb{R}$ be a continuous function. Without loss of generality, assume $f(a) < f(b)$. We choose an arbitrary $c \in \mathbb{R}$ such that $f(a) \leq c \leq f(b)$ and define the set:
    $$X \coloneqq \{t \in [a, b] \mid f(t) \leq c\}$$
    Observe that $X \neq \emptyset$ (since $a \in X$ as $f(a) \leq c$) and $X$ is bounded above by $b$. Thus, by the completeness axiom of the real numbers, the supremum $s \coloneqq \sup X$ exists, and clearly $s \in [a, b]$.

    By the definition of the supremum, for all $n \in \mathbb{N}$, there exists $t_n \in X \cap [s - \frac{1}{n}, s]$. 
    Because $t_n \to s$ as $n \to \infty$ and $f$ is continuous, we have $\lim_{n\to\infty} f(t_n) = f(s)$. 
    Since $t_n \in X$, we know $f(t_n) \leq c$ for all $n \in \mathbb{N}$. By the limit-preserving properties of inequalities, it follows that $f(s) \leq c$.

    Now, assume for the sake of contradiction that $f(s) < c$. 
    Since $f(s) < c \leq f(b)$, we know $s \neq b$, which implies $s < b$. 

    Let $\epsilon = c - f(s) > 0$. By the continuity of $f$ at $s$, there exists a $\delta > 0$ such that for all $y \in [a, b]$:
    $$|y - s| < \delta \implies |f(y) - f(s)| < \epsilon$$
    This implies $f(y) < f(s) + \epsilon = c$. 

    Now, define $y \coloneqq \min(b, s + \frac{\delta}{2})$. 
    Since $s < b$ and $\delta > 0$, we strictly have $y > s$. Furthermore, $y \in [a, b]$ and $|y - s| = \frac{\delta}{2} < \delta$. 
    Therefore, by our continuity implication, $f(y) < c$, which means $y \in X$. 

    However, we have found an element $y \in X$ such that $y > s$, which directly contradicts the fact that $s$ is an upper bound (the supremum) of $X$. 

    Thus, our assumption that $f(s) < c$ must be false. Since we have already established that $f(s) \leq c$, we conclude that $f(s) = c$.
\end{proof}


\begin{definition}
A bounded and closed interval is called \textbf{compact}.
\end{definition}

\begin{example}
    The interval $[0, 1]$ is compact while $[0, \infty)$ is not.
\end{example}

\begin{lemma}
    Let $[a, b]$ be a compact interval and let $(x_n)_{n\in\mathbb{N}} \subset [a, b]$ be a sequence. Then there exists a subsequence $(x_{n_k})_{k\in\mathbb{N}}$ such that 
    $$\lim_{k\to\infty} x_{n_k} = x_0 \text{ for some } x_0 \in [a, b].$$
\end{lemma}

\begin{definition}
    Let $f: D \subset \mathbb{R} \to \mathbb{R}$ be a function
    \begin{itemize}
        \item If for $x_0 \in D$ it holds that $f(x_0) \geq f(x) \forall x \in D$, then $f$ has a local maximum at $x_0$.
        \item If for $x_0 \in D$ it holds that $f(x_0) \leq f(x) \forall x \in D$, then $f$ has a local minimum at $x_0$.
        \item Local maxima and minima are called local extrema.
    \end{itemize}
\end{definition}

\begin{theorem}
    Let $f: I \subset \mathbb{R} \to \mathbb{R}$ be a continous function and $I$ be a compact interval. Then $f$ is bounded and attains its maximum and minimum value on $I$. Hence
    $$\exists x_{min}, x_{max} \in I \text{ such that } f(x_{min}) \leq f(x) \leq f(x_{max}) \forall x \in I$$
\end{theorem}

